\hypertarget{tag-000C}{}\label{app:rederive:fam-src-A.5.1}
\begin{proof}
For the standard twist data, the classification of
\cite[Theorem~1.4, p.~3]{DuEtAl2026v1}, with the exact parameter
identification of \cref{night:P5.7}, is re-derived in
\cref{night:P5.6}. It gives $X_c\simeq X_{c'}$ exactly when
$c-c'\in2\mathbb Z$. The explicit maps on the regular, finite,
and cusp charts are those of \cref{fam:src:A.3.1}, displayed in
\eqref{wb:pshift:eq:full-isomorphism} and satisfying
\eqref{wb:pshift:eq:coherent-shift-action}.

As a complex parameter space for the isomorphism classes represented
by this marked family, the quotient $\mathcal P/(2\mathbb Z)$ has
the explicit coordinate
\begin{equation}
 w=\exp(-\pi ic),\qquad 0<|w|<\exp(-\pi M).
 \hypertarget{tag-03CL}{}\label{wb:pshift:eq:parameter-disc}
\end{equation}
Indeed $|w|=\exp(\pi\operatorname{Im}c)$, and two complex numbers
have the same value of this exponential exactly when they differ
by an even integer. Every point in the displayed punctured disc
has a logarithm giving a $c\in\mathcal P$. The exponential has
nonzero derivative, so these local inverse logarithms give the
claimed quotient complex structure. The integer-period quotient
seen by the middle family instead has coordinate $w^2$; it is
covered two-to-one by \eqref{wb:pshift:eq:parameter-disc}. Its two lifts are distinguished by
the exact order-two summand in \eqref{wb:pshift:eq:finite-class}. The resulting parameter space classifies the members of the marked family.
Proposition~\ref{fam:src:A.3.1} proves properness of this quotient
family and its joint holomorphic dependence. Its Kodaira--Spencer
class is nonzero by Proposition~\ref{fam:src:A.4.2}.
\end{proof}

